\documentclass[a4,11pt]{aleph-notas}
% Se puede ver la documentación aquí:
% https://github.com/alephsub0/LaTeX_aleph-notas

% -- Paquetes adicionales
\usepackage{enumitem}
\usepackage{url}
\usepackage{array}
\usepackage{booktabs}
\usepackage{longtable}
\hypersetup{
    urlcolor=blue,
    linkcolor=blue,
}

% -- Datos
\institucion{Facultad de Ciencias Exactas, Naturales y Ambientales}
\carrera{Catálogo STEM}
\asignatura{Álgebra lineal}
\tema{Tema 06: Espacios vectoriales y subespacios vectoriales}
\autor{Andrés Merino}
\fecha{Periodo 2026-01}

\logouno[0.16\textwidth]{Logos/logoPUCE_04_ac}
\definecolor{colortext}{HTML}{1957d1}
\definecolor{colordef}{HTML}{1957d1}
\fuente{montserrat}

\begin{document}

\encabezado 

%%%%%%%%%%%%%%%%%%%%%%%%%%%%%%%%%%%%%%%%
\section*{Resultado de Aprendizaje}
%%%%%%%%%%%%%%%%%%%%%%%%%%%%%%%%%%%%%%%%

%%%%%%%%%%%%%%%%%%%%%%%%%%%%%%%%%%%%%%%%
\subsection*{RdA de la asignatura:}
\begin{itemize}[leftmargin=*]
    \item \textbf{RdA 1:} Comprender los conceptos fundamentales del Álgebra Lineal, incluyendo el estudio de matrices, determinantes, sistemas de ecuaciones lineales y espacios vectoriales, destacando su importancia en el análisis y resolución de problemas matemáticos.
    \item \textbf{RdA 2:} Resolver problemas prácticos y teóricos relacionados con el Álgebra Lineal, utilizando técnicas de cálculo con matrices, determinantes y sistemas de ecuaciones, así como en el análisis de espacios vectoriales.    
    % \item \textbf{RdA 3:} Aplicar los principios y métodos del Álgebra Lineal en una variedad de contextos demostrando comprensión de su relevancia y utilidad en situaciones prácticas y teóricas.
\end{itemize}


%%%%%%%%%%%%%%%%%%%%%%%%%%%%%%%%%%%%%%%%
\subsection*{RdA de la actividad:}
\begin{itemize}[leftmargin=*]
    \item Comprender los axiomas que definen un espacio vectorial.
    \item Aplicar el criterio de subespacio para determinar si un subconjunto de un espacio vectorial es un subespacio. 
\end{itemize}

%%%%%%%%%%%%%%%%%%%%%%%%%%%%%%%%%%%%%%%%
\section*{Introducción}
%%%%%%%%%%%%%%%%%%%%%%%%%%%%%%%%%%%%%%%%

%%%%%%%%%%%%%%%%%%%%%%%%%%%%%%%%%%%%%%%%
\paragraph{Control de lectura:}
Antes de iniciar la clase, se aplicará el control de lectura del siguiente enlace:
\href{https://gemini.google.com/share/cb5551a4eb1e}{Control de lectura - Clase 06}.

\paragraph{Pregunta inicial:}
¿Podemos extender las ideas de suma y producto en $\mathbb{R}^n$ a otros conjuntos?

%%%%%%%%%%%%%%%%%%%%%%%%%%%%%%%%%%%%%%%%
\section*{Desarrollo}
%%%%%%%%%%%%%%%%%%%%%%%%%%%%%%%%%%%%%%%%

%%%%%%%%%%%%%%%%%%%%%%%%%%%%%%%%%%%%%%%%
\subsection*{Actividad 1: Axiomas de espacio vectorial y criterio de subespacio}

La clase inicia con una recuperación breve de operaciones en $\mathbb{R}^n$ y del papel del campo de escalares. Luego, mediante clase magistral y práctica guiada, se formalizan los axiomas de espacio vectorial, se analizan ejemplos y contraejemplos, y se aplica el criterio de subespacio en distintos contextos.

\paragraph{¿Cómo lo haremos?}
\begin{itemize}[leftmargin=*]
    \item \textbf{Clase magistral:} se presentan definición y axiomas de espacio vectorial, propiedades básicas, y ejemplos en $\mathbb{R}^n$, matrices, funciones y polinomios. Se usa el resumen \href{https://fcena-puce.github.io/AlgebraLineal-05-N0048/01Resumen/Resumen06.pdf}{Resumen06.pdf}.
    \item \textbf{Resolución de ejercicios:} se guiará a los estudiantes en la solución de ejercicios del documento \href{https://fcena-puce.github.io/AlgebraLineal-05-N0048/04Ejercicios/Ejercicios06.pdf}{Ejercicios06.pdf}, verificando qué axiomas se cumplen o fallan en operaciones no usuales.
\end{itemize}

\paragraph{Verificación de aprendizaje:}
\begin{enumerate}[leftmargin=*]
    \item ¿Qué axiomas permiten concluir que un conjunto con operaciones dadas es un espacio vectorial?
    \item ¿Cómo se aplica el criterio de subespacio para decidir si $W\subseteq E$ es subespacio de $E$?
\end{enumerate}

%%%%%%%%%%%%%%%%%%%%%%%%%%%%%%%%%%%%%%%%
\section*{Cierre}
%%%%%%%%%%%%%%%%%%%%%%%%%%%%%%%%%%%%%%%%

\paragraph{Tarea:}
Resolver del libro \href{https://catalogobiblioteca.puce.edu.ec/cgi-bin/koha/opac-detail.pl?biblionumber=86081}{Fundamentos de álgebra lineal de Larson}, sección 4.2, los ejercicios: 13, 15, 17,25, 26; sección 4.3, los ejercicios: 3, 5, 7, 11, 13, 16.

\paragraph{Pregunta de investigación:}
\begin{enumerate}[leftmargin=*]
    \item ¿Existen otros conjuntos con operaciones distintas a las usuales que formen espacios vectoriales en el campo $\mathbb{R}$?
    \item ¿Qué es el campo de Galois y cómo se relaciona con los espacios vectoriales?
    \item ¿Las palabras pueden formar un espacio vectorial? 
\end{enumerate}

\paragraph{Para la próxima clase:}  
Visualizar el video \href{https://www.youtube.com/watch?v=k7RM-ot2NWY\&list=PLZHQObOWTQDPD3MizzM2xVFitgF8hE\_ab\&index=2}{Combinaciones lineales, espacio generado y vectores base}.
Realizar la actividad de clase invertida: \href{https://fcena-puce.github.io/AlgebraLineal-05-N0048/07Act-ClaseInvertida/02ClsInv-Independencia.pdf}{02ClsInv-Independencia.pdf}.


\end{document}
